\section*{NeurIPS Paper Checklist}

\begin{enumerate}

\item {\bf Claims}
    \item[] Question: Do the main claims made in the abstract and introduction accurately reflect the paper's contributions and scope?
    \item[] Answer: \answerYes{}
    \item[] Justification: The four contributions stated in Section~1 (helical PE, set-prediction formulation with Hungarian matching, $\sim$\,$420\times$ end-to-end speed-up at higher accuracy, cross-cell transfer) each map directly to a results subsection (Section~4) and an ablation row (Table~\ref{tab:ablate}).

\item {\bf Limitations}
    \item[] Question: Does the paper discuss the limitations of the work performed by the authors?
    \item[] Answer: \answerYes{}
    \item[] Justification: Section~6 explicitly discusses panel scope ($7$ TFs, $2$ cell types), saturation of $F_1$, fixed slot count $K$ and clipping behavior, dependence on peak-call supervision, and limited novel-motif discovery.

\item {\bf Theory assumptions and proofs}
    \item[] Question: For each theoretical result, does the paper provide the full set of assumptions and a complete (and correct) proof?
    \item[] Answer: \answerNA{}
    \item[] Justification: The paper makes empirical, not theoretical, claims. No theorems are stated.

\item {\bf Experimental result reproducibility}
    \item[] Question: Does the paper fully disclose all the information needed to reproduce the main experimental results of the paper to the extent that it affects the main claims and/or conclusions of the paper (regardless of whether the code and data are provided or not)?
    \item[] Answer: \answerYes{}
    \item[] Justification: Architecture, hyperparameters, optimizer, schedule, data splits, peak-calling settings, and evaluation protocol are described in Sections~3--4 and Appendix~A. Anonymized code and pretrained weights are included in the supplementary ZIP.

\item {\bf Open access to data and code}
    \item[] Question: Does the paper provide open access to the data and code, with sufficient instructions to faithfully reproduce the main experimental results, as described in supplemental material?
    \item[] Answer: \answerYes{}
    \item[] Justification: All data are public ENCODE ChIP-seq and JASPAR PWMs (cited and licensed). The supplementary ZIP contains training/evaluation scripts, environment file, pretrained weights for the K562-7tf model, and a README with exact commands.

\item {\bf Experimental setting/details}
    \item[] Question: Does the paper specify all the training and test details (e.g., data splits, hyperparameters, how they were chosen, type of optimizer) necessary to understand the results?
    \item[] Answer: \answerYes{}
    \item[] Justification: Section~3.3 and Appendix~A specify optimizer, LR schedule, loss weights, batch size, epochs, and chromosome-level data splits.

\item {\bf Experiment statistical significance}
    \item[] Question: Does the paper report error bars suitably and correctly defined or other appropriate information about the statistical significance of the experiments?
    \item[] Answer: \answerYes{}
    \item[] Justification: Table~\ref{tab:main} reports mean $\pm$ standard deviation over $5$ random seeds (variation is over weight initialization and slot init noise; data splits are fixed). Ablation deltas in Table~\ref{tab:ablate} use the same $5$-seed protocol; only mean values are shown to save space.

\item {\bf Experiments compute resources}
    \item[] Question: For each experiment, does the paper provide sufficient information on the computer resources (type of compute workers, memory, time of execution) needed to reproduce the experiments?
    \item[] Answer: \answerYes{}
    \item[] Justification: Appendix~A reports per-run cost ($\sim$\,$10$ A100-hours), total reported-experiment cost ($\sim$\,$300$ A100-hours), and an estimate of preliminary/failed runs ($\sim$\,$300$ additional A100-hours).

\item {\bf Code of ethics}
    \item[] Question: Does the research conducted in the paper conform, in every respect, with the NeurIPS Code of Ethics?
    \item[] Answer: \answerYes{}
    \item[] Justification: All data are public, de-identified ENCODE assays. No human subjects, no scraped data, no dual-use risks beyond those general to genomics modeling.

\item {\bf Broader impacts}
    \item[] Question: Does the paper discuss both potential positive societal impacts and negative societal impacts of the work performed?
    \item[] Answer: \answerYes{}
    \item[] Justification: Section~7 discusses positive impacts (faster, lower-cost binding catalogues; reproducibility) and the absence of specific identified misuse risks beyond standard genomics-model considerations.

\item {\bf Safeguards}
    \item[] Question: Does the paper describe safeguards that have been put in place for responsible release of data or models that have a high risk for misuse?
    \item[] Answer: \answerNA{}
    \item[] Justification: The released model predicts TF binding from public DNA sequences and does not generate content, identify individuals, or have a recognized high-risk misuse pathway.

\item {\bf Licenses for existing assets}
    \item[] Question: Are the creators or original owners of assets used in the paper properly credited and are the license and terms of use explicitly mentioned and properly respected?
    \item[] Answer: \answerYes{}
    \item[] Justification: ENCODE data are credited and used under the ENCODE Consortium's data-use terms; JASPAR 2024 PWMs are used under CC0; PyTorch (BSD), MEME Suite (academic), and TF-MoDISco (MIT) are credited and licensed for the reported uses. Specific URLs and versions are listed in the supplementary README.

\item {\bf New assets}
    \item[] Question: Are new assets introduced in the paper well documented and is the documentation provided alongside the assets?
    \item[] Answer: \answerYes{}
    \item[] Justification: The released BEACON model, code, and pretrained K562-7tf weights are documented with a model card and README in the anonymized supplementary ZIP, including training data, intended use, and known limitations.

\item {\bf Crowdsourcing and research with human subjects}
    \item[] Question: For crowdsourcing experiments and research with human subjects, does the paper include the full text of instructions given to participants and screenshots, if applicable, as well as details about compensation (if any)?
    \item[] Answer: \answerNA{}
    \item[] Justification: No crowdsourcing or human-subjects research is conducted.

\item {\bf Institutional review board (IRB) approvals or equivalent for research with human subjects}
    \item[] Question: Does the paper describe potential risks incurred by study participants, whether such risks were disclosed to the subjects, and whether IRB approvals (or an equivalent approval/review based on the requirements of your country or institution) were obtained?
    \item[] Answer: \answerNA{}
    \item[] Justification: No human-subjects research; ENCODE data were generated and released by the ENCODE Consortium under their own ethical and consent procedures.

\item {\bf Declaration of LLM usage}
    \item[] Question: Does the paper describe the usage of LLMs if it is an important, original, or non-standard component of the core methods in this research?
    \item[] Answer: \answerNA{}
    \item[] Justification: LLMs are not part of the core methodology. Authors used general-purpose writing assistants only for grammar and formatting, which the policy does not require declaring.

\end{enumerate}
